\section{Fundamentação teórica}

\subsection{Arquiteturas de LLM relacionadas ao problema}

Os componentes mais diretamente relacionados a este trabalho são, de um lado, os LLMs generativos baseados em transformers \parencite{vaswani2017attention} e, de outro, os modelos de embeddings utilizados para recuperação semântica \parencite{nie2024textEmbeddingSurvey}. Enquanto o LLM atua como componente de síntese textual, os embeddings convertem documentos e consultas em vetores, permitindo o cálculo de similaridade semântica em bancos vetoriais \parencite{gao2023ragSurvey}.

Em sistemas RAG, esses componentes operam de maneira complementar: o modelo de embedding sustenta a recuperação, enquanto o modelo generativo formula a resposta a partir do contexto recuperado \parencite{lewisEtAl2020rag,gao2023ragSurvey}. Essa separação favorece uma arquitetura modular, a especialização de componentes e a evolução incremental da solução \parencite{gao2024modularRag}.

Segundo \textcite{gao2023ragSurvey}, sistemas RAG podem ser classificados em abordagens ingênuas, avançadas e modulares. Para o problema investigado, a arquitetura modular mostra-se adequada por permitir substituir, de forma independente, extratores documentais, estratégias de segmentação, modelos de embedding, bancos vetoriais, reranqueadores e modelos gerativos. Em contextos institucionais, essa característica é desejável porque infraestrutura, custos, segurança e políticas internas podem exigir alterações graduais da pilha tecnológica.

Essa modularidade, contudo, não deve ser interpretada como garantia automática de qualidade. A literatura apresenta o RAG como estratégia para reduzir problemas de desatualização, baixa verificabilidade e alucinação, mas não como solução suficiente para eliminá-los. O desempenho final continua dependente da qualidade documental, da segmentação, da recuperação, do desenho do prompt, da seleção de contexto e dos mecanismos de avaliação adotados \parencite{gao2023ragSurvey,gupta2024comprehensiveRag,saadFalconEtAl2024ares}. Essa leitura é importante em aplicações públicas, nas quais uma resposta plausível, mas mal fundamentada, pode comprometer confiança institucional.

\subsection{Técnicas relacionadas}

As principais técnicas relacionadas ao \textit{pipeline} deste projeto incluem extração documental, \textit{chunking}, geração de \textit{embeddings} semânticos, recuperação vetorial, seleção contextual e \textit{prompting} orientado a evidências, podendo também envolver etapas de reranqueamento em arquiteturas RAG \parencite{gao2023ragSurvey,maharjan2026chunkingReranking,wangEtAl2025segmentation}. Em conjunto, essas técnicas convertem o acervo em unidades recuperáveis, identificam trechos semanticamente relevantes e delimitam o contexto encaminhado ao modelo generativo. O \textit{prompting} orientado a evidências complementa esse fluxo ao instruir o modelo a privilegiar o contexto recuperado, evitar inferências especulativas e explicitar referências verificáveis.

Dentre essas técnicas, o \textit{chunking} ocupa posição central. Estudos sobre segmentação documental em RAG indicam que escolhas inadequadas de divisão podem comprometer a recuperação semântica, seja por fragmentarem indevidamente o contexto, seja por gerarem blocos excessivamente extensos, com perda de precisão na busca \parencite{wangEtAl2025segmentation,bhatEtAl2025chunkSize}. No presente projeto, a segmentação documental exigiu sucessivos ajustes empíricos para equilibrar granularidade semântica, preservação de contexto e restrições operacionais da base vetorial e do processo de ingestão. Assim, a experiência prática corrobora a literatura ao indicar que o \textit{chunking} não constitui etapa meramente operacional, mas um dos principais determinantes do desempenho de um sistema RAG.

\textit{Surveys} e \textit{frameworks} também indicam que, em ambientes reais, a qualidade de sistemas RAG depende de mecanismos de avaliação, rastreabilidade e comparação entre configurações, não apenas do \textit{pipeline} básico de recuperação e geração \parencite{gao2023ragSurvey,saadFalconEtAl2024ares,esEtAl2023ragas}. Essa perspectiva orientou a criação de \textit{scripts} próprios de preparação, importação, avaliação e documentação da configuração experimental.

No campo da avaliação, estratégias baseadas em LLM como juiz têm sido utilizadas para escalar a análise de respostas geradas, sobretudo quando combinadas com rubricas explícitas \parencite{liuEtAl2023geval,liuEtAl2025judgeRag}. Ainda assim, o julgamento automatizado pode variar conforme o modelo avaliador adotado, a rubrica utilizada e o grau de abertura da tarefa; por isso, em aplicações RAG, tende a ser mais robusto quando tratado como instrumento auxiliar, combinado a evidências qualitativas, validação humana e análise das fontes recuperadas \parencite{liuEtAl2025judgeRag,liuEtAl2023geval,saadFalconEtAl2024ares,esEtAl2023ragas}. Benchmarks voltados ao domínio jurídico reforçam essa preocupação ao tratar a avaliação de RAG como problema que envolve não apenas qualidade textual da resposta, mas também recuperação de documentos pertinentes, fidelidade ao contexto e adequação da evidência utilizada \parencite{pipitoneAlami2024legalbenchRag}.

\subsection{Trabalhos e soluções relacionadas}

\textcite{gao2023ragSurvey} oferecem uma base conceitual abrangente para compreender o RAG como resposta aos problemas de desatualização, baixa verificabilidade e geração não rastreável em LLMs, além de orientar a decomposição do problema em etapas de recuperação, aumento de contexto e geração. No plano arquitetural, \textcite{gupta2024comprehensiveRag} consolidam o estado da arte em RAG ao discutir como mecanismos de recuperação e geração se combinam para mitigar limitações de LLMs e ao destacar desafios de precisão, fidelidade e coordenação entre componentes.

Os trabalhos de \textcite{minaee2024llmSurvey} e \textcite{vaswani2017attention} oferecem o pano de fundo técnico para o uso de LLMs generativos neste projeto: o primeiro sintetiza a evolução e as capacidades desses modelos em ampla gama de tarefas, enquanto o segundo estabelece a base arquitetural dos transformers. \textcite{nie2024textEmbeddingSurvey} complementam esse quadro ao tratar os \textit{text embeddings} como tecnologia fundacional para tarefas de busca semântica e recuperação de informação, enfatizando a relevância da integração entre LLMs e \textit{embeddings} em aplicações contemporâneas.

A literatura sobre alucinação em LLMs evidencia que esses modelos podem produzir saídas fluentemente redigidas, porém factualmente incorretas ou não verificáveis, o que reforça, em contextos públicos sensíveis, a necessidade de soluções explicáveis e auditáveis, com mecanismos de fundamentação documental, governança e controle institucional \parencite{anhHoang2025hallucinations,gao2023ragSurvey,deAlmeidaSantos2025aiGovernance}. Em domínios jurídicos, esse risco é particularmente relevante porque respostas plausíveis podem induzir interpretações normativas, citações inexistentes ou inferências de autoridade sem base documental suficiente \parencite{dahlEtAl2024largeLegalFictions}.

Em termos de contexto institucional, \textcite{un2015responsiveGovernance}, \textcite{ipu2024aiParliaments}, \textcite{wfd2023guidelinesAiParliaments} e \textcite{geunis2023parliamentaryMonitoring} indicam que parlamentos e administrações públicas operam sob pressões crescentes por transparência, responsividade, responsabilização e uso qualificado da informação. Nesse domínio, discursos e debates constituem insumos centrais para análise política e apoio técnico, mas os sítios parlamentares frequentemente funcionam mais como repositórios informacionais do que como instrumentos de exploração substantiva, tornando o acesso dependente de conhecimento prévio sobre estrutura e vocabulário \parencite{skubicFiser2024discourseReview,bandeiraBernardes2016information}.

No caso brasileiro, os Textos para Discussão da Consultoria Legislativa do Senado oferecem evidência empírica sobre a relevância e a complexidade do acervo de pronunciamentos. \textcite{blanco2025plenarioPalanqueEstudio} analisa discursos em plenário entre 2007 e 2024 e identifica mudanças de volume, extensão, tempo de fala e interatividade, associadas a transformações institucionais e tecnológicas. Em estudo complementar, \textcite{blanco2026letrasGraudas} examina leiturabilidade e traços morfossintáticos do mesmo domínio discursivo por meio de técnicas computacionais de análise textual. Esses trabalhos reforçam que os discursos do Senado constituem corpus relevante para métodos computacionais, embora se concentrem em análise descritiva e interpretativa, e não no desenvolvimento de uma camada interativa voltada à consulta, à recuperação semântica e à geração de respostas verificáveis.

No campo da governança pública de inteligência artificial, estudos e diretrizes institucionais enfatizam responsabilidade, transparência, monitoramento, gestão de riscos e alinhamento a finalidades públicas \parencite{deAlmeidaSantos2025aiGovernance,ipu2024aiParliaments,wfd2023guidelinesAiParliaments}. Para aplicações RAG em acervos legislativos, esse conjunto de referências reforça a importância de bases documentais controladas, metadados preservados, respostas verificáveis e avaliação contínua.

No domínio político, \textcite{khaliq2024ragar} demonstram que arquiteturas RAG podem apoiar tarefas de fact-checking político com raciocínio fundamentado em evidências externas, articulando recuperação rigorosa e geração apoiada em fontes verificáveis. Ainda que o objetivo deste trabalho não seja a verificação factual multimodal, a pesquisa desses autores é relevante por demonstrar a aderência do domínio político-parlamentar a abordagens baseadas em recuperação e geração com evidências.

Por fim, \textcite{martelloteViola2025organizacaoConhecimento} aproxima a discussão do contexto brasileiro ao tratar da inteligência artificial na organização do conhecimento legislativo. Sua contribuição reforça a pertinência de soluções explicáveis, documentalmente ancoradas e informacionalmente estruturadas em ambientes parlamentares, oferecendo referencial específico para a aplicação proposta neste trabalho.

Trabalhos recentes tornam mais visível a aproximação entre RAG, LLMs, conhecimento legislativo e acesso a informação parlamentar. A Tabela~\ref{tab:trabalhos-relacionados} sintetiza os estudos mais próximos identificados na revisão e explicita sua relação com a presente proposta. Em conjunto, eles mostram que há aplicações relevantes em fact-checking político, redação legislativa, recuperação de legislação, interfaces conversacionais e debates parlamentares \parencite{khaliq2024ragar,chouhanGertz2024lexdrafter,rogiersEtAl2024kamerraad,matoshiEtAl2025parliamentRag,garrisonEtAl2024talkNdaa,colombo2024legisAi,colomboEtAl2025legisSearch,mosbachEtAl2025worldAvatarParliament,sarnikar2025legislativeTexts}. Ainda assim, diferem do presente trabalho por privilegiarem outros corpus, línguas, jurisdições, tarefas ou estratégias de avaliação.

\begin{scriptsize}
\begin{longtable}{@{}p{0.17\textwidth}p{0.22\textwidth}p{0.22\textwidth}p{0.28\textwidth}@{}}
\caption{Trabalhos relacionados mais próximos e posicionamento da proposta}
\label{tab:trabalhos-relacionados}\\
\hline
\textbf{Trabalho} & \textbf{Corpus ou domínio} & \textbf{Abordagem e avaliação} & \textbf{Relação com este artigo} \\
\hline
\endfirsthead
\hline
\textbf{Trabalho} & \textbf{Corpus ou domínio} & \textbf{Abordagem e avaliação} & \textbf{Relação com este artigo} \\
\hline
\endhead
\textcite{blanco2025plenarioPalanqueEstudio} e \textcite{blanco2026letrasGraudas} & Discursos em plenário do Senado Federal brasileiro. & Análises computacionais descritivas sobre volume, extensão, interatividade, leiturabilidade e complexidade textual. & Compartilham o domínio empírico brasileiro, mas não propõem sistema RAG, interface conversacional ou avaliação de respostas geradas. \\
\textcite{martelloteViola2025organizacaoConhecimento} & Dados e informações legislativas do Parlamento brasileiro. & Discussão aplicada sobre IA e organização do conhecimento legislativo. & Aproxima o problema ao contexto brasileiro de organização da informação, mas não avalia um chatbot RAG sobre discursos parlamentares. \\
\textcite{khaliq2024ragar} & Fact-checking político multimodal. & Arquitetura RAG para raciocínio apoiado em evidências externas. & Demonstra a utilidade de recuperação com evidências no domínio político, embora com foco em verificação factual, e não em consulta a acervo legislativo. \\
\textcite{chouhanGertz2024lexdrafter} & Documentos legislativos da União Europeia, com foco em definições terminológicas. & RAG para apoiar redação de artigos de definição; código disponibilizado pelos autores. & É diretamente legislativo e usa RAG, mas a tarefa principal é redação terminológica, não consulta auditável a discursos parlamentares. \\
\textcite{rogiersEtAl2024kamerraad} & Política nacional belga. & Sumarização hierárquica e interface conversacional para recuperação de informação política. & Aproxima-se pela interface conversacional em domínio parlamentar, mas difere por corpus, língua, desenho de sumarização e foco avaliativo. \\
\textcite{matoshiEtAl2025parliamentRag} & Acervos e negócios parlamentares suíços. & Projetos-piloto com modelagem de tópicos e RAG para melhorar acesso digital. & É uma das aproximações mais diretas ao problema de acesso parlamentar, mas em outra jurisdição e com foco em indexação temática institucional. \\
\textcite{garrisonEtAl2024talkNdaa} & Legislação congressional dos Estados Unidos, em especial o NDAA. & Avaliação de RAG para sumarização e consulta de legislação longa. & Contribui para comparação em legislação e avaliação de RAG, mas não trata discursos parlamentares nem o contexto brasileiro. \\
\textcite{colomboEtAl2025legisSearch} & Legislação italiana representada em grafo. & Busca legislativa com grafo, LLM, embeddings, comparação com BM25/TF-IDF e ablação. & É metodologicamente forte para recuperação legislativa, mas prioriza navegação em leis e relações normativas, não discursos e metadados parlamentares. \\
\textcite{mosbachEtAl2025worldAvatarParliament} & Debates parlamentares do Bundestag alemão. & Sistema RAG híbrido com grafo de conhecimento e recuperação vetorial. & Aproxima-se por tratar debates parlamentares, mas usa arquitetura híbrida com grafo e outro contexto linguístico-institucional. \\
\textcite{sarnikar2025legislativeTexts} & Textos legislativos longos. & Uso de LLMs para consulta e compreensão de legislação extensa. & Dialoga com a dificuldade de consultar textos legislativos longos, mas não enfatiza corpus de discursos, metadados parlamentares ou prova de conceito reprodutível. \\
\hline
\end{longtable}
\end{scriptsize}

\subsection{Síntese da revisão e lacuna de pesquisa}

A revisão realizada permite identificar três conjuntos de literatura complementares. O primeiro, voltado a RAG, LLMs, \textit{embeddings}, segmentação documental e avaliação automatizada, oferece a base técnica do trabalho \parencite{lewisEtAl2020rag,gao2023ragSurvey,minaee2024llmSurvey,nie2024textEmbeddingSurvey,saadFalconEtAl2024ares}. O segundo, dedicado a parlamentos, organização da informação legislativa e discurso parlamentar, delimita a relevância institucional do domínio \parencite{skubicFiser2024discourseReview,martelloteViola2025organizacaoConhecimento,ipu2024aiParliaments}. O terceiro, mais diretamente ligado ao caso brasileiro, aplica métodos computacionais ao corpus de discursos do Senado, mas com foco predominante em análise descritiva, linguística ou discursiva \parencite{blanco2025plenarioPalanqueEstudio,blanco2026letrasGraudas}. 

A lacuna remanescente, portanto, não está na ausência de estudos sobre RAG nem na ausência de estudos sobre discursos parlamentares, mas na articulação aplicada entre esses dois campos. Embora a literatura indique o potencial do RAG para elevar a precisão factual e a verificabilidade de sistemas baseados em LLM, esse potencial depende de decisões metodológicas situadas, e não apenas da adoção abstrata da arquitetura \parencite{gao2023ragSurvey,gupta2024comprehensiveRag,saadFalconEtAl2024ares}. Parte significativa da literatura técnica permanece concentrada em discussões conceituais amplas, cenários corporativos ou aplicações em domínios distintos do acervo legislativo brasileiro \parencite{gao2023ragSurvey,gupta2024comprehensiveRag}. Mesmo os trabalhos mais próximos em legislação, política e parlamentos enfatizam fact-checking, redação de normas, recuperação de legislação, grafos de conhecimento ou outras jurisdições \parencite{khaliq2024ragar,chouhanGertz2024lexdrafter,rogiersEtAl2024kamerraad,matoshiEtAl2025parliamentRag,garrisonEtAl2024talkNdaa,colomboEtAl2025legisSearch,mosbachEtAl2025worldAvatarParliament,sarnikar2025legislativeTexts}. Ainda que pesquisas já explorem computacionalmente discursos do Senado sob perspectivas quantitativas, discursivas e linguísticas, são menos frequentes estudos que articulem, em um mesmo desenho aplicado, corpus parlamentar nacional, preparação reprodutível de base documental, preservação de metadados legislativos, fluxo RAG operacional e avaliação empírica das respostas geradas \parencite{blanco2025plenarioPalanqueEstudio,blanco2026letrasGraudas,martelloteViola2025organizacaoConhecimento}.

Assim, identifica-se uma lacuna aplicada relacionada à construção, à documentação e à avaliação de soluções RAG voltadas especificamente à consulta de discursos parlamentares no contexto do Senado Federal. Este trabalho busca contribuir para seu preenchimento ao propor, implementar e avaliar uma prova de conceito reprodutível de \textit{chatbot} legislativo baseada em RAG, com foco em discursos da 56\textordfeminine{} Legislatura. A contribuição reside não apenas no protótipo, mas também na explicitação do \textit{pipeline}, no enriquecimento dos \textit{chunks} com metadados, na documentação dos artefatos de ingestão e avaliação e na combinação entre estudos de caso, rodadas automatizadas com rubrica e análise manual amostral.
